\documentclass[11pt]{article}

\usepackage[margin=1in]{geometry}
\usepackage[T1]{fontenc}
\usepackage[utf8]{inputenc}
\usepackage{lmodern}
\usepackage{microtype}
\usepackage{amsmath,amssymb,amsthm,mathtools}
\usepackage{booktabs}
\usepackage{array}
\usepackage{longtable}
\usepackage{tabularx}
\usepackage{hyperref}
\usepackage[nameinlink,noabbrev]{cleveref}
\usepackage{enumitem}
\usepackage{graphicx}
\usepackage{xcolor}

\hypersetup{
  colorlinks=true,
  linkcolor=blue,
  citecolor=blue,
  urlcolor=blue,
  pdftitle={A Compact-Fiber Structural Prediction of the Low-Energy Electromagnetic Coupling},
  pdfauthor={A. R. Wells}
}

\title{A Compact-Fiber Structural Prediction of the Low-Energy Electromagnetic Coupling}

\author{A.\ R.\ Wells\\
Dual-Frame Research Group}

\date{\today}

\begin{document}

\maketitle

\begin{abstract}
We derive a compact-fiber structural prediction for a low-energy
electromagnetic bridge coupling. The compact-fiber carrier
\(F=Z_4\times Z_4\times Z_8\) supplies \(|F|=128\), and the
radiative/vibrational bridge-placement readout supplies
\(|V_{\rm rad}^{\rm bridge}|=3^2=9\), giving the base bridge count
\(B=137\). Cover-mediated feedback gives \(K=74/15\), and the first
closed return yields
\[
x^4-137x^3-\frac{74}{15}x^2+\frac{296}{105}=0.
\]
The positive root is
\[
x_{\rm fiber}=137.0359991771859\ldots,
\qquad
\alpha_F=1/x_{\rm fiber}.
\]
No measured value of \(\alpha\) is used as an algebraic input, and no
continuously adjusted coefficient appears in the bridge equation. The
identification \(\alpha_F=\alpha_{\rm low-energy}\) is made by physical
role: leading atomic binding reads the material/radiative bridge through
two vertices and therefore carries an \(\alpha_F^2\) dependence. The
numerical comparison with the empirical fine-structure constant is a
consistency check, not an error-budget closure; route-specific
metrological extraction remains downstream.
\end{abstract}

\tableofcontents

\section{Introduction}

The fine-structure constant is one of the central dimensionless constants
of low-energy physics. In conventional notation,
\[
\alpha=\frac{e^2}{4\pi\varepsilon_0\hbar c},
\]
or, in naturalized electromagnetic units, it is the dimensionless
coupling governing the strength of electromagnetic interactions. It
appears in atomic binding, spectroscopy, radiative corrections,
scattering amplitudes, and precision metrology. Its inverse value is
approximately
\[
\alpha^{-1}\approx 137.035999177.
\]

In standard physics this value is not derived from a deeper structural
principle. It is inferred through operational routes such as atom recoil,
spectroscopy, electron \(g-2\), and electrical standards, and then used
as an input to low-energy quantum electrodynamics and atomic theory
\cite{NIST_CODATA_2022_alpha_inv,Parker2018_alpha_Cs,
Morel2020_alpha_Rb,Aoyama2019_electron_gminus2_review,
BIPM_SI_Brochure_9th}.

The purpose of this paper is to present a compact-fiber structural
prediction for the low-energy electromagnetic bridge coupling. The
derived object is denoted
\[
\alpha_F,
\qquad
x=\alpha_F^{-1}.
\]
The central result is that compact-fiber carrier structure,
radiative/vibrational bridge placement, cover-mediated feedback,
self-consistent transport sharing, and the first closed return yield a
quartic equation for \(x\). Its positive root is
\[
x_{\rm fiber}=137.0359991771859\ldots,
\]
so that
\[
\alpha_F=\frac{1}{x_{\rm fiber}}.
\]
The measured value of \(\alpha\) is not used as an algebraic input to
the bridge equation.

\subsection{Structural target}

The compact-fiber framework begins from the finite carrier
\[
F=Z_4\times Z_4\times Z_8,
\]
with
\[
|F|=128,
\qquad
N_m=4,
\qquad
N_e=8.
\]
The two \(Z_4\) factors represent the paired magnetic sector, while the
\(Z_8\) factor represents the electric-cover sector. The relation
\[
N_e=2N_m
\]
supplies the canonical electric-cover / magnetic-base quotient used in
the cover-mediated feedback term. The present paper does not rederive
the compact fiber from first principles; it takes the carrier and sector
capacities from the foundational compact-fiber construction and derives
the electromagnetic bridge coupling from them
\cite{Wells2026PaperI,Wells2026PaperII,Wells2026PaperIII}.

The electromagnetic coupling is treated here as a carrier-level
material/radiative bridge normalization. It is not an element-by-element
function of specific labels
\[
f\in F,
\qquad
v\in V_{\rm rad}.
\]
Instead,
\[
\alpha_F(f,v)=\alpha_F
\]
for all bridge-admissible element labels. The relevant structural
question is therefore how material carrier access and
radiative/vibrational bridge-placement access combine at the level of
the electromagnetic bridge, not how many ordered pairs exist in
\[
F\times V_{\rm rad}.
\]

\subsection{Derivation strategy}

The derivation proceeds through four structural steps. First, the
material compact-fiber carrier and the radiative/vibrational
bridge-placement carrier determine the base bridge count
\[
B=128+9=137.
\]
Second, cover-mediated feedback from the paired magnetic sector into the
electromagnetic bridge determines the feedback burden
\[
K=\frac{74}{15}.
\]
Third, self-consistent transport sharing gives the first-order
fixed-point structure
\[
x=B+\frac{K}{x}.
\]
Fourth, the first closed return of the admitted feedback supplies the
second-order return term. These ingredients yield the compact bridge
equation
\[
x
=
B+\frac{K}{x}
-
\frac{K N_m}{(N_e-1)x^3}.
\]
The detailed derivation of each term is given in the sections that
follow.

\subsection{Identification with the low-energy electromagnetic coupling}

The fixed-point derivation produces a compact-fiber structural bridge
coupling
\[
\alpha_F=\frac{1}{x_{\rm fiber}}.
\]
The identification
\[
\alpha_F=\alpha_{\rm low-energy}
\]
is made by physical role, not by numerical agreement alone. The
compact-fiber electromagnetic bridge is the material/radiative interface.
A leading atomic binding observable is a closed material-radiative
source-response process with two bridge vertices,
\[
\text{material source}
\rightarrow
\text{radiative/electromagnetic bridge}
\rightarrow
\text{material response}.
\]
Thus the leading atomic scale carries two bridge factors,
\[
E_{\rm atomic,F}
=
C_{\rm mass/readout}\alpha_F^2.
\]
This matches the role of \(\alpha^2\) in the Hartree/Rydberg scale,
where
\[
E_h=\alpha^2m_ec^2,
\qquad
R_\infty=\frac{\alpha^2m_ec}{2h}.
\]
This paper therefore treats \(\alpha_F\) as the compact-fiber structural
counterpart of the low-energy fine-structure constant.

\subsection{Scope of the paper}

The paper separates the structural coupling derivation from operational
measurement routes. It derives the compact-fiber bridge coupling
\(\alpha_F\), identifies it with the low-energy electromagnetic coupling
by its \(\alpha_F^2\) role in leading atomic binding, and gives a
controlled correspondence to empirical extraction routes.

The result is a conditional structural derivation: given the named
compact-fiber framework lemmas, the fixed-point equations and their roots
follow algebraically. It is not presented as an unconditional derivation
from bare postulates alone. The paper also does not claim a completed
route-by-route derivation of CODATA or metrological extraction methods;
full route-specific metrological closure remains downstream work.

The measured value of \(\alpha\) is used only after the derivation, as
an empirical consistency check. No measured value of \(\alpha\) is used
to set \(B\), \(K\), the closed-return coefficient, or any continuously
adjusted coefficient in the bridge equation.
\section{Compact-Fiber Carrier and Electromagnetic Bridge}

\subsection{Compact-fiber carrier}

The compact-fiber carrier is
\[
F=Z_4\times Z_4\times Z_8.
\]
This carrier and its sector structure are established in the foundational compact-fiber work and used in subsequent atomic, representational, and radiation-sector applications \cite{Wells2026PaperI,Wells2026PaperII,Wells2026PaperIII,Wells2026Radiation}.

Its cardinality is
\[
|F|=4\cdot4\cdot8=128.
\]

The two \(Z_4\) factors supply the paired magnetic sector. The \(Z_8\) factor supplies the electric-cover sector. The corresponding sector capacities are
\[
N_m=4,
\qquad
N_e=8.
\]

The relation
\[
N_m\mid N_e,
\qquad
8=2\cdot4,
\]
gives the canonical electric-cover / magnetic-base quotient
\[
\pi:Z_8\to Z_4.
\]
This quotient is not independent bridge data. It is downstream of the compact-fiber sector capacities. In the alpha derivation, \(\pi\) becomes relevant when paired magnetic feedback is read through the electric-cover structure.

\subsection{Material/radiative bridge}

The electromagnetic coupling is treated as a carrier-level bridge normalization between material closure and radiative/vibrational placement.

Material closure is represented by
\[
F=Z_4\times Z_4\times Z_8.
\]
Radiative/vibrational placement is represented by a bridge-placement carrier
\[
V_{\rm rad}^{\rm bridge}.
\]
The electromagnetic bridge is therefore the material/radiative interface:
\[
F
\longleftrightarrow
V_{\rm rad}^{\rm bridge}.
\]

The coupling derived here is denoted
\[
\alpha_F.
\]
Its inverse is
\[
x=\alpha_F^{-1}.
\]
The quantity \(x\) is the effective compact-fiber electromagnetic bridge capacity.

\subsection{Carrier-level universality}

Within the compact-fiber framework, the electromagnetic bridge is a carrier-to-carrier material/radiative interface. Its coupling is therefore universal at the bridge level. For admissible element labels
\[
f\in F,
\qquad
v\in V_{\rm rad}^{\rm bridge},
\]
the bridge coupling satisfies
\[
\alpha_F(f,v)=\alpha_F.
\]

Thus the electromagnetic coupling is not an element-level quantity whose value depends on a selected compact-fiber state and a selected radiative-placement state. It is a scalar per-channel bridge share determined by carrier-level access capacity.

This carrier-level universality is the reason the alpha bridge is not normalized over the pair-configuration space
\[
F\times V_{\rm rad}^{\rm bridge}.
\]
A pair space would be appropriate for an observable whose elementary alternatives are ordered pairs \((f,v)\). The alpha bridge has a different role: it is the scalar normalization of the material/radiative interface itself.

The relevant structural inputs are therefore carrier-level invariants:
\[
|F|,
\qquad
|V_{\rm rad}^{\rm bridge}|,
\qquad
N_m,
\qquad
N_e.
\]
For the compact fiber,
\[
|F|=128,
\qquad
N_m=4,
\qquad
N_e=8.
\]
The remaining carrier input is the radiative/vibrational bridge-placement count
\[
|V_{\rm rad}^{\rm bridge}|,
\]
which is derived below.

\subsection{\texorpdfstring{Why the bridge is not counted as \(F\times V_{\rm rad}^{\rm bridge}\)}{Why the bridge is not counted as F x Vrad bridge}}

A Cartesian-product count
\[
|F\times V_{\rm rad}^{\rm bridge}|
=
|F|\,|V_{\rm rad}^{\rm bridge}|
\]
would be appropriate if an electromagnetic bridge state required simultaneous specification of one material element and one radiative-placement element as a joint state label.

That is not the role of the electromagnetic bridge in this derivation.

The bridge coupling is not an element-pair table. It is a scalar carrier-level bridge normalization. The material carrier and the radiative/vibrational bridge-placement carrier supply mechanism-distinct access to the bridge. They are independently sufficient carrier-level access structures, not simultaneously required labels of a joint microscopic state.

Therefore the base bridge count uses disjoint-union access,
\[
F\sqcup V_{\rm rad}^{\rm bridge},
\]
not product-state access,
\[
F\times V_{\rm rad}^{\rm bridge}.
\]
Thus the base count is
\[
B=|F|+|V_{\rm rad}^{\rm bridge}|.
\]

This distinction is load-bearing. It separates the compact-fiber alpha derivation from a pair-counting argument.

\subsection{\texorpdfstring{Structural role of \(x=\alpha_F^{-1}\)}{Structural role of x = alphaF inverse}}

The quantity
\[
x=\alpha_F^{-1}
\]
is not a primitive compact-fiber cardinality. It is the self-consistent effective bridge capacity obtained after the base carrier count is corrected by cover-mediated feedback.

The derivation proceeds in two stages. First, the base bridge count
\[
B
\]
is obtained from carrier-level material and radiative/vibrational bridge-placement access. Second, the net cover-mediated feedback burden
\[
K
\]
is shared over the final effective bridge capacity
\[
x.
\]

This gives the first-order fixed point
\[
x=B+\frac{K}{x}.
\]
A closed-return correction then refines this to
\[
x
=
B+\frac{K}{x}
-
\frac{K N_m}{(N_e-1)x^3}.
\]

The positive root of the resulting quartic is the compact-fiber structural prediction for
\[
\alpha_F^{-1}.
\]

\subsection{Structural output}

The output of this derivation is the compact-fiber electromagnetic bridge coupling
\[
\alpha_F.
\]
The later identification
\[
\alpha_F=\alpha_{\rm low-energy}
\]
is made by physical role, not by using the measured fine-structure constant as an input.
\section{Base Bridge Count}

\subsection{Material carrier contribution}

The material side of the electromagnetic bridge is the compact fiber
\[
F=Z_4\times Z_4\times Z_8.
\]
Therefore the material carrier contributes
\[
|F|=4\cdot4\cdot8=128.
\]
This count is fixed once the compact fiber is fixed.

\subsection{Radiative/vibrational bridge-placement contribution}

The radiative/vibrational bridge-placement carrier contributes
\[
|V_{\rm rad}^{\rm bridge}|=3^2=9.
\]
This count is not a magnetic-sector count and is not inserted to force
the number \(137\). Its provenance is radiative/vibrational bridge
placement.

A one-dimensional scalar-displacement relation with a distinguished datum
has the qualitative sign-status quotient
\[
\{-,0,+\},
\]
corresponding to position below, at, or above the datum. The displacement
algebra itself is not reduced to three elements; rather, the
pre-projection placement readout has three sign-status classes.

Radiation is treated here as a one-dimensional vibratory seed. In
three-dimensional placement geometry, the orientation space of a
one-dimensional axis is two-dimensional:
\[
\dim \operatorname{Gr}(1,3)=2.
\]
Thus the electromagnetic bridge reads the radiative/vibrational placement
over two bridge-facing placement dimensions. The radiation-sector paper
treats this as the vibrational placement layer of the radiation
architecture \cite{Wells2026Radiation}.

Given the two-dimensional placement readout and the ternary sign-status
quotient in each placement dimension,
\[
V_{\rm rad}^{\rm bridge}
\cong
\{-,0,+\}^2,
\]
and therefore
\[
|V_{\rm rad}^{\rm bridge}|
=
|\{-,0,+\}^2|
=
3^2
=
9.
\]
Only the placement count is used in the alpha derivation. No group law on
\(V_{\rm rad}^{\rm bridge}\) is required.

\subsection{Disjoint-union access count}

The base bridge count uses the carrier-access composition rule:
\[
\boxed{
\text{mechanism-distinct, independently sufficient access routes add;}
}
\]
\[
\boxed{
\text{simultaneously required labels multiply;}
}
\]
\[
\boxed{
\text{mere refinements or marginals of the same mechanism collapse rather than add.}
}
\]

This rule separates three cases. If two labels must be specified
simultaneously to identify a single elementary event, the correct access
space is a Cartesian product. If two structures supply different-kind,
independently sufficient access routes to the same scalar bridge
observable, the correct access space is a disjoint union. If one
structure is only a refinement or marginalization of the other, it is not
counted as an additional independent access route.

The electromagnetic bridge is a carrier-level material/radiative
interface. It is not a joint state space whose elements are ordered pairs
\[
(f,v)\in F\times V_{\rm rad}^{\rm bridge}.
\]
A product count would give
\[
|F\times V_{\rm rad}^{\rm bridge}|
=
|F|\,|V_{\rm rad}^{\rm bridge}|.
\]
That would be the correct count only if an electromagnetic bridge event
required simultaneous specification of one material element and one
radiative-placement element as a joint state label.

The bridge coupling derived here is instead a scalar carrier-level
normalization. Material carrier access and radiative/vibrational
bridge-placement access are mechanism-distinct and independently
sufficient at the bridge level. They are not mutually indexed degrees of
freedom of a single joint state, and neither is merely a marginal
refinement of the other.

Thus the relevant access space is the disjoint union
\[
F\sqcup V_{\rm rad}^{\rm bridge},
\]
not the tensor/product-state access space
\[
F\times V_{\rm rad}^{\rm bridge}.
\]
The base bridge count is therefore
\[
B
=
|F\sqcup V_{\rm rad}^{\rm bridge}|
=
|F|+|V_{\rm rad}^{\rm bridge}|.
\]
Substituting the two carrier counts gives
\[
B=128+9=137.
\]

\subsection{Base count result}

The base electromagnetic bridge count is
\[
\boxed{
B=137.
}
\]
This is the zeroth-order compact-fiber bridge capacity before
cover-mediated feedback is included. It is not the final value of
\[
x=\alpha_F^{-1}.
\]
The final effective bridge capacity is obtained by adding feedback
self-consistently,
\[
x=B+\frac{K}{x},
\]
and then including the closed-return correction
\[
-\frac{K N_m}{(N_e-1)x^3}.
\]

The dependency structure is:
\[
F=Z_4\times Z_4\times Z_8
\quad\Rightarrow\quad
|F|=128,
\]
\[
\{-,0,+\}\ \text{over two radiative placement dimensions}
\quad\Rightarrow\quad
|V_{\rm rad}^{\rm bridge}|=3^2=9,
\]
\[
\text{mechanism-distinct disjoint-union access}
\quad\Rightarrow\quad
B=|F|+|V_{\rm rad}^{\rm bridge}|.
\]
Thus
\[
\boxed{
B=137
}
\]
is a structural base count determined by carrier access and
bridge-placement readout, not a fitted numerical input.
\section{Cover-Mediated Feedback Burden}

\subsection{Feedback burden from cover-mediated readout}

The base bridge count
\[
B=137
\]
is not the final effective bridge capacity. The compact fiber also
supplies cover-mediated feedback from the paired magnetic sector into
the electromagnetic bridge.

The paired magnetic sector is
\[
Z_4\times Z_4,
\]
and the electric-cover sector is
\[
Z_8.
\]
Since
\[
N_m\mid N_e,
\qquad
8=2\cdot4,
\]
the compact fiber contains the canonical electric-cover / magnetic-base
quotient
\[
\pi:Z_8\to Z_4.
\]
This quotient is derived from the sector-capacity relation and is not
independent bridge data.

The target observable
\[
\alpha_F
\]
is a carrier-level electromagnetic bridge coupling. Paired magnetic
feedback therefore does not enter as a purely magnetic-internal quantity;
it is read through the electric-cover / magnetic-base interface.

\subsection{Canonical cover-interface readout routes}

The base-cover feedback channel has two mechanism-distinct scalar
readout resolutions. The first is direct electric-cover resolution,
\[
\frac{1}{N_e}.
\]
The second is fully resolved joint magnetic/electric resolution,
\[
\frac{1}{N_mN_e}.
\]
These are not sequential filters applied to the same event. Sequential
transport would multiply capacities; mechanism-distinct scalar access
contributions to a single feedback burden add before the burden is
distributed over the final bridge capacity.

Thus the gross cover-mediated readout is
\[
\frac1{N_e}+\frac1{N_mN_e}.
\]
At common denominator \(N_mN_e\),
\[
\frac1{N_e}+\frac1{N_mN_e}
=
\frac{N_m}{N_mN_e}
+
\frac1{N_mN_e}
=
\frac{N_m+1}{N_mN_e}.
\]
The numerator
\[
N_m+1
\]
is the intrinsic mechanism count of the base-cover feedback channel.
For
\[
N_m=4,
\]
this gives
\[
N_m+1=5.
\]
Therefore the gross feedback burden is
\[
\boxed{
M_{\rm fb}=5.
}
\]

The local product-fiber denominator
\[
N_mN_e
\]
is appropriate when the target observable is resolved at product-fiber
scale. That is not the normalization used for the alpha bridge equation.
The target observable is the effective bridge capacity
\[
x=\alpha_F^{-1},
\]
so the bridge share is
\[
\frac1x.
\]
Reusing the local denominator \(N_mN_e\) and then dividing again by
\(x\) would double-count target normalization. For the alpha bridge, the
feedback contribution has the form
\[
\frac{K}{x},
\]
not
\[
\frac{K}{N_mN_e\,x}.
\]
Thus the intrinsic feedback numerator \(N_m+1\) is carried forward as
the gross feedback burden.

\subsection{Scalar self-resolution cost}

The paired magnetic contrast space is
\[
G=Z_4\times Z_4,
\]
with
\[
|G|=N_m^2=16.
\]
The datum state
\[
(0,0)
\]
carries no paired magnetic contrast and is excluded from the active
contrast support. Define
\[
G^\ast=G\setminus\{(0,0)\}.
\]
Then
\[
|G^\ast|=N_m^2-1.
\]
For \(N_m=4\),
\[
|G^\ast|=16-1=15.
\]

The electromagnetic bridge coupling is a scalar carrier-level observable.
It does not resolve the individual labels of \(G^\ast\). The unresolved
non-datum contrast support therefore contributes one scalar
self-resolution cost distributed over the active contrast space:
\[
A_{\rm self}
=
\frac1{|G^\ast|}
=
\frac1{N_m^2-1}.
\]
For \(N_m=4\),
\[
A_{\rm self}
=
\frac1{15}.
\]
This is a self-resolution cost, not an additional feedback source and
not an attenuation factor applied independently to every feedback event.

\subsection{Closure-budget accounting}

The closure budget contains the gross generated feedback burden
\[
M_{\rm fb}=N_m+1
\]
and the consumed self-resolution cost
\[
A_{\rm self}=\frac1{N_m^2-1}.
\]
The source-level budget relation is
\[
\text{gross generated feedback}
=
\text{delivered feedback}
+
\text{consumed self-resolution cost}.
\]
Equivalently,
\[
\text{delivered feedback}
=
\text{gross generated feedback}
-
\text{consumed self-resolution cost}.
\]
Thus
\[
K=M_{\rm fb}-A_{\rm self}.
\]

This excludes a multiplicative form such as
\[
M_{\rm fb}(1-A_{\rm self})
\]
in the present role assignment. A multiplicative form would treat
\(A_{\rm self}\) as an attenuation applied to each feedback event. Here
\(A_{\rm self}\) is an additive cost consumed from the same closure
budget before the net burden is delivered to the bridge.

Substituting the compact-fiber values,
\[
K
=
(N_m+1)
-
\frac1{N_m^2-1}.
\]
For \(N_m=4\),
\[
K
=
5-\frac1{15}
=
\frac{74}{15}.
\]
Therefore the net first-order feedback burden is
\[
\boxed{
K=\frac{74}{15}.
}
\]

\subsection{\texorpdfstring{Role of \(K\) in the bridge equation}{Role of K in the bridge equation}}

The quantity
\[
K=\frac{74}{15}
\]
is not a direct channel count to be added to \(B\). It is a feedback
burden delivered into the effective bridge capacity. Therefore it appears
as
\[
\frac{K}{x}
\]
in the self-consistent bridge equation, not as
\[
K.
\]
The first-order equation is
\[
x=B+\frac{K}{x},
\]
not
\[
x=B+K,
\]
and not merely a one-pass correction
\[
x=B+\frac{K}{B}.
\]
The denominator is the final effective bridge capacity \(x\), because
the feedback modifies the same bridge capacity over which it is
distributed.

The feedback-burden dependency structure is therefore:
\[
\text{cover-mediated gross burden}
\quad
N_m+1=5,
\]
\[
\text{active non-datum contrast support}
\quad
N_m^2-1=15,
\]
\[
\text{scalar self-resolution cost}
\quad
\frac1{N_m^2-1}=\frac1{15},
\]
\[
\text{net delivered feedback}
\quad
K=(N_m+1)-\frac1{N_m^2-1}.
\]
Thus
\[
\boxed{
K=\frac{74}{15}
}
\]
is a structural feedback burden determined by the compact-fiber closure
budget, not a fitted constant.
\section{Self-Consistent Fixed Point}

\subsection{Effective bridge capacity}

The previous sections established the base electromagnetic bridge count
\[
B=137
\]
and the net first-order feedback burden
\[
K=\frac{74}{15}.
\]
The inverse compact-fiber electromagnetic bridge coupling is denoted
\[
x=\alpha_F^{-1}.
\]
The quantity \(x\) is not a primitive carrier cardinality. The primitive material carrier has
\[
|F|=128,
\]
and the base bridge count is
\[
B=128+9=137.
\]
The effective bridge capacity \(x\) is obtained only after cover-mediated feedback is admitted and distributed over the same bridge capacity that it modifies.

The central point is that \(K\) is not a new carrier pool and not a direct count to be added to \(B\). It is a feedback burden delivered into the electromagnetic bridge. The appropriate first-order equation is therefore self-consistent:
\[
x=B+\frac{K}{x}.
\]

\subsection{\texorpdfstring{Why the feedback term is \(K/x\)}{Why the feedback term is K/x}}

If the feedback burden \(K\) were an independent carrier pool, the first-order equation would be
\[
x=B+K.
\]
That is not its role. The burden is delivered into the existing bridge and distributed over the effective bridge capacity. A burden distributed over \(N\) effective channels contributes as
\[
\frac{K}{N}.
\]

If the bridge capacity were fixed at its uncorrected base value \(B\), the one-pass perturbative correction would be
\[
x\approx B+\frac{K}{B}.
\]
But the feedback changes the bridge capacity that appears in the denominator. The denominator must therefore be the final effective capacity
\[
x,
\]
not the uncorrected base count
\[
B.
\]
Thus the first-order feedback contribution is
\[
\frac{K}{x}.
\]

\subsection{First-order bridge equation}

The first-order compact-fiber bridge equation is
\[
\boxed{
x=B+\frac{K}{x}.
}
\]
Substituting
\[
B=137,
\qquad
K=\frac{74}{15},
\]
gives
\[
x=137+\frac{74}{15x}.
\]
Multiplying by \(x\),
\[
x^2=137x+\frac{74}{15}.
\]
Therefore
\[
\boxed{
x^2-137x-\frac{74}{15}=0.
}
\]
This is the first-order fixed-point equation for the effective electromagnetic bridge capacity.

\subsection{Positive root}

The positive root is
\[
x_{\rm quad}
=
\frac{137+\sqrt{137^2+4(74/15)}}{2}.
\]
Since
\[
4\left(\frac{74}{15}\right)
=
\frac{296}{15},
\]
this may also be written as
\[
x_{\rm quad}
=
\frac{137+\sqrt{137^2+296/15}}{2}.
\]
Numerically,
\[
\boxed{
x_{\rm quad}=137.03600027236\ldots.
}
\]
The corresponding first-order compact-fiber bridge coupling is
\[
\alpha_{F,\rm quad}
=
\frac{1}{x_{\rm quad}}.
\]

\subsection{First-order accuracy and residual scale}

The first-order result is structurally nontrivial. It fixes the inverse electromagnetic bridge coupling close to the observed low-energy inverse fine-structure constant scale without using the measured value of \(\alpha\) as an input.

It is not, however, the final compact-fiber prediction. At modern precision, the first-order value remains high by approximately
\[
8\text{ ppb}
\]
relative to the quoted low-energy inverse fine-structure constant. The remaining discrepancy motivates the closed-return analysis of the admitted first-order feedback burden.

\subsection{Independence from the closed-return correction}

The first-order equation
\[
x=B+\frac{K}{x}
\]
depends only on the self-consistent transport-sharing rule: a feedback burden delivered into an effective bridge capacity is distributed over the final capacity that it helps determine.

Under that rule, the first-order result follows algebraically. It does not depend on the second-order closed-return correction. If the closed-return term were revised, the first-order compact-fiber bridge lemma would remain intact.
\section{Closed-Return Correction}

\subsection{Closed-return structure}

The first-order bridge equation is
\[
x=B+\frac{K}{x}.
\]
Its positive root is
\[
x_{\rm quad}=137.03600027236\ldots.
\]
This first-order result is structurally constrained but not sufficient at
modern comparison precision. The admitted cover-mediated feedback has a
first closed return on the scalar two-node bridge graph. That return
supplies the second-order correction
\[
\boxed{
\Delta x_{\rm 2nd}
=
-\frac{K N_m}{(N_e-1)x^3}.
}
\]

\subsection{The two-node feedback graph}

The relevant feedback graph contains two transport nodes:
\[
\text{EM bridge}
\]
and
\[
\text{paired magnetic source}.
\]
The cover map
\[
\pi:Z_8\to Z_4
\]
is channel-internal, not an independent third transport node. A structure
would qualify as an independent transport node only if it had independent
state space, independent transport capacity, sequential multiplicative
factorization, and independent normalization.

The cover map \(\pi\) fails these criteria. It is a morphism, not a
carrier object with its own independent position space. Its kernel is
internal to \(Z_8\), and elements in a fiber are co-accessed by
base-level operations factoring through \(\pi\). Additive cover-to-base
readout therefore has the signature of mechanism-distinct routes through
a channel, not sequential traversal through a separate transport node.

Thus the active scalar feedback graph is
\[
\text{EM bridge}
\longleftrightarrow
\text{paired magnetic source}.
\]

\subsection{First closed return}

The first-order term
\[
+\frac{K}{x}
\]
is an outward cover-mediated feedback contribution from the paired
magnetic source into the electromagnetic bridge. The first closed return
is the shortest circuit that leaves the bridge and returns to it through
the paired magnetic source:
\[
\text{EM bridge}
\rightarrow
\text{paired magnetic source}
\rightarrow
\text{EM bridge}.
\]
This return circuit requires two additional bridge traversals beyond the
first-order feedback term. Each traversal is diluted over the effective
bridge capacity \(x\), so the two additional traversals contribute
\[
\frac1{x^2}.
\]
The return therefore has the scale
\[
\frac{K}{x}\cdot\frac1{x^2}
=
\frac{K}{x^3}.
\]

\subsection{Return coefficient}

The return coefficient is
\[
\frac{N_m}{N_e-1}.
\]
The numerator \(N_m\) is the magnetic reception capacity of the paired
magnetic side. The denominator \(N_e-1\) is the electric-cover complement
capacity remaining after the first-order source activation.

This denominator follows from co-access plus transport accounting. A
base-level activation through
\[
\pi:Z_8\to Z_4
\]
co-accesses the full fiber
\[
\pi^{-1}(m),
\]
but it does so as one carrier-level transport event, not as independent
electric activations for each element of the fiber. Therefore the
first-order event removes one transport-event unit from the electric-cover
transport capacity, leaving
\[
N_e-1.
\]
It does not remove \(d\) independent electric sites, so the denominator
is not \(N_e-d\). It is also not \(N_e\), because the first-order source
activation has already committed one electric-cover transport event
before the return is formed.

For the compact fiber,
\[
N_m=4,
\qquad
N_e=8,
\]
and therefore
\[
\frac{N_m}{N_e-1}
=
\frac47.
\]

\subsection{Negative sign}

The sign of the closed return is fixed by its role as derivative
self-response of the first-order feedback term. The first-order feedback
contribution is
\[
g(x)=\frac{K}{x}.
\]
Since
\[
K>0,
\qquad
x>0,
\]
we have
\[
g'(x)=-\frac{K}{x^2}<0.
\]
Thus, if the closed return is the derivative self-response of the
first-order feedback to the bridge capacity that feedback modifies, it
must reduce the outward first-order contribution. The path factors
\[
\frac1{x^2}
\]
and
\[
\frac{N_m}{N_e-1}
\]
are positive. They affect the magnitude but not the sign. Therefore
\[
\Delta x_{\rm 2nd}<0.
\]

\subsection{Second-order correction}

Combining the first-order admitted feedback, the two additional bridge
traversals, the return coefficient, and the negative derivative
self-response sign gives
\[
\Delta x_{\rm 2nd}
=
-\frac{K}{x}
\cdot
\frac1{x^2}
\cdot
\frac{N_m}{N_e-1}.
\]
Therefore
\[
\boxed{
\Delta x_{\rm 2nd}
=
-\frac{K N_m}{(N_e-1)x^3}.
}
\]

For the compact-fiber values,
\[
K=\frac{74}{15},
\qquad
N_m=4,
\qquad
N_e=8,
\]
the coefficient is
\[
\frac{K N_m}{N_e-1}
=
\frac{(74/15)\cdot4}{7}
=
\frac{296}{105}.
\]
Thus the second-order correction is
\[
\boxed{
\Delta x_{\rm 2nd}
=
-\frac{296}{105x^3}.
}
\]

The measured fine-structure constant is not used to set
\[
\frac47
\]
or any other coefficient in the corrected fixed-point equation.

\subsection{Corrected fixed point}

Adding the closed-return correction to the first-order equation gives
\[
x
=
B+\frac{K}{x}
-
\frac{K N_m}{(N_e-1)x^3}.
\]
Substituting
\[
B=137,
\qquad
K=\frac{74}{15},
\qquad
N_m=4,
\qquad
N_e=8,
\]
we obtain
\[
x
=
137+\frac{74}{15x}
-
\frac{296}{105x^3}.
\]
This is the second-order compact-fiber bridge equation.

\subsection{Status of the correction}

The admitted second-order correction is conditionally unique within the
current scalar two-node bridge graph. At order \(O(x^{-3})\), possible
sources would include a second independent first-order feedback source,
a new carrier pool, a direct cubic mechanism, base-feedback interference,
a non-scalar correction, or higher-order leakage. Under the current scalar
two-node graph, none of these alternatives is admitted: the closed return
of the already-admitted first-order feedback is the only same-order scalar
closed-return correction.

This is a conditional uniqueness claim. It is relative to the current
scalar bridge graph and the stated transport-node, source/complement,
derivative self-response, and scalar-readout rules. Future radiation-side
feedback mechanisms would require re-audit, but no such mechanism is
admitted in the present scalar bridge graph.

If repeated closed returns are admitted, each additional round trip is
suppressed by at least
\[
\frac{C}{x^2},
\qquad
C=\frac{N_m}{N_e-1}=\frac47.
\]
For
\[
x\approx137.036,
\]
this suppression factor is approximately
\[
3.0\times10^{-5}.
\]
Thus higher-order repeated returns, if present, are operationally small
at the current comparison scale. A complete all-orders return theorem
remains open; the present paper uses only the first closed-return
correction.
\section{Numerical Prediction}

\subsection{From fixed point to polynomial}

The previous section derived the second-order compact-fiber bridge
equation
\[
x
=
137+\frac{74}{15x}
-
\frac{296}{105x^3}.
\]
Multiplying by \(x^3\) converts this fixed-point equation into a quartic
polynomial whose positive root defines the compact-fiber inverse bridge
coupling:
\[
\boxed{
x^4
-
137x^3
-
\frac{74}{15}x^2
+
\frac{296}{105}
=
0.
}
\]
The measured value of \(\alpha\) is not used to obtain the equation or
its root.

\subsection{Positive root}

Solving
\[
x^4
-
137x^3
-
\frac{74}{15}x^2
+
\frac{296}{105}
=
0
\]
gives the physically relevant positive root
\[
\boxed{
x_{\rm fiber}=137.0359991771859\ldots.
}
\]
Therefore the compact-fiber electromagnetic bridge coupling is
\[
\boxed{
\alpha_F
=
\frac1{x_{\rm fiber}}.
}
\]
This is the central numerical output of the compact-fiber bridge
equation.

\subsection{First-order and second-order values}

The first-order equation
\[
x=B+\frac{K}{x},
\qquad
B=137,
\qquad
K=\frac{74}{15},
\]
gives
\[
\boxed{
x_{\rm quad}=137.03600027236\ldots.
}
\]
The second-order closed-return equation replaces \(x_{\rm quad}\) with
the positive quartic root \(x_{\rm fiber}\). The shift is produced by
the admitted scalar two-node return term
\[
-\frac{296}{105x^3}.
\]
The first-order result fixes the inverse bridge value close to the
empirical low-energy inverse fine-structure constant scale but remains
high by approximately
\[
8.0\ \text{ppb}.
\]
The second-order closed return supplies the refinement admitted by the
framework rules stated above.

\subsection{Comparison with the empirical low-energy value}

Using the CODATA 2022 recommended value
\cite{NIST_CODATA_2022_alpha_inv,NIST_CODATA_2022_wallet}
\[
\alpha^{-1}_{\rm CODATA}=137.035999177(21),
\]
the compact-fiber value
\[
x_{\rm fiber}=137.0359991771859\ldots
\]
lies within the experimental uncertainty band of the quoted value. The
difference from the CODATA 2022 central value is approximately
\[
1.4\times10^{-3}\ \text{ppb},
\]
while the quoted experimental uncertainty is approximately
\[
0.15\ \text{ppb}.
\]

The numerical comparison is not used to construct the bridge equation.
It is also not a statistical validation of the framework, since the
present paper does not assign a theoretical uncertainty to the imported
bridge-readout rules, conditional framework lemmas, or possible omitted
same-order mechanisms. The comparison is therefore reported as a
consistency check between the compact-fiber point value and the empirical
low-energy coupling, not as an error-budget closure.

\subsection{Status of the prediction}

No measured value of \(\alpha\) appears in the compact-fiber bridge
equation, and no continuously adjusted coefficient appears in the
polynomial. The coefficients are fixed structurally before comparison
with CODATA:
\[
137=128+9,
\qquad
\frac{74}{15}
=
5-\frac1{15},
\qquad
\frac{296}{105}
=
\left(\frac{74}{15}\right)\left(\frac47\right).
\]
Their structural origins are
\[
128=|Z_4\times Z_4\times Z_8|,
\qquad
9=|V_{\rm rad}^{\rm bridge}|=3^2,
\]
\[
5=N_m+1,
\qquad
\frac1{15}=\frac1{N_m^2-1},
\qquad
\frac47=\frac{N_m}{N_e-1}.
\]

Given the stated compact-fiber bridge lemmas, the quartic equation and
its positive root follow algebraically. The result remains conditional
on the bridge-readout, closure-budget, and closed-return rules stated
above; the paper does not claim that those rules are derived here from
bare postulates alone. The role identification
\[
\alpha_F=\alpha_{\rm low-energy}
\]
is supported by the numerical agreement, but is not based on numerical
agreement alone. Its structural basis is the material/radiative bridge
role of \(\alpha_F\) and the two-vertex \(\alpha_F^2\) dependence of
leading atomic binding.
\section{Low-Energy Electromagnetic Role Identification}

The preceding sections derive a compact-fiber structural bridge coupling
\[
\alpha_F=\frac1{x_{\rm fiber}},
\]
where \(x_{\rm fiber}\) is the positive root of the compact-fiber bridge polynomial. The remaining question is the physical role of this structural coupling. The identification
\[
\alpha_F=\alpha_{\rm low-energy}
\]
is not made by numerical agreement alone. It is made by role: the compact-fiber electromagnetic bridge is the material/radiative interface, and leading atomic binding reads that bridge through a two-vertex source-response process.

\subsection{Electromagnetic bridge as material/radiative interface}

In the compact-fiber framework, material closure is represented by
\[
F=Z_4\times Z_4\times Z_8,
\]
while radiative/vibrational placement is represented by
\[
V_{\rm rad}.
\]

The electromagnetic interaction is identified as the carrier-level bridge between these two structural regimes:
\[
F
\longleftrightarrow
V_{\rm rad}.
\]

The coupling associated with this bridge is
\[
\alpha_F.
\]

It is not an element-level label depending on a particular pair
\[
(f,v)\in F\times V_{\rm rad}.
\]
It is the scalar bridge normalization of the material/radiative interface itself. This is why the derivation solves for
\[
x=\alpha_F^{-1}
\]
as an effective bridge capacity.

\subsection{Two-vertex source-response structure}

A leading atomic binding observable is not a one-way bridge event. It is a closed source-response process:
\[
\text{material source}
\rightarrow
\text{radiative/electromagnetic bridge}
\rightarrow
\text{material response}.
\]

The first bridge vertex carries material structure into the electromagnetic/radiative bridge. The second bridge vertex returns the electromagnetic response to the material sector. Thus the leading atomic source-response scale contains two bridge factors:
\[
\alpha_F\cdot\alpha_F=\alpha_F^2.
\]

Therefore the compact-fiber leading atomic scale has the form
\[
E_{\rm atomic,F}
=
C_{\rm mass/readout}\alpha_F^2,
\]
where \(C_{\rm mass/readout}\) contains dimensional mass, action, and finite-readout matching factors. The finite-readout and matching-correction layer is treated separately in \cite{Wells2026FiniteReadout}. The load-bearing structural fact is the coupling power:
\[
E_{\rm atomic,F}\propto \alpha_F^2.
\]

\subsection{Correspondence with Hartree/Rydberg scaling}

In conventional low-energy atomic physics, the Hartree energy is
\[
E_h=\alpha^2m_ec^2.
\]
The Rydberg constant is
\[
R_\infty=\frac{\alpha^2m_ec}{2h}.
\]

Both expressions show the same structural role:
\[
\text{atomic binding scale}\propto\alpha^2.
\]

The compact-fiber source-response structure gives
\[
E_{\rm atomic,F}
=
C_{\rm mass/readout}\alpha_F^2.
\]

Thus the compact-fiber bridge coupling occupies the same low-energy role as the conventional fine-structure constant:
\[
\alpha_F
\leftrightarrow
\alpha_{\rm low-energy}.
\]

This is the physical role identification used in this paper.

\subsection{Boundary of the role identification}

The role identification claims that \(\alpha_F\) is the compact-fiber structural counterpart of the low-energy electromagnetic coupling.

It does not independently derive the dimensional anchors or correction pipelines that appear in particular empirical routes. Those are route-specific readout problems.

The alpha theorem derived in this paper supplies the structural electromagnetic bridge coupling. Other constants and correction layers determine how this coupling appears in specific measurement pipelines.

\subsection{Relation to the Planck/action bridge}

Some empirical routes to \(\alpha\), especially Rydberg/recoil and
\(h/m\)-based routes, involve the action anchor
\[
h.
\]

The compact-fiber program treats this as a separate structural problem.
The Planck/action and gravitational-coupling analysis is developed in a
separate constants-sector paper \cite{Wells2026PlanckGravity}. That
work gives the Planck/action bridge its own structural content,
schematically of the form
\[
h=\frac{m_n c S_n^2}{R_{\rm eff}},
\]
where \(R_{\rm eff}\) is a dimensionless compact-fiber structural factor
and the remaining quantities are dimensional anchors connecting the
structural/natural-unit system to conventional units.

The Planck/action bridge was already identified as a structural scaling
problem in the radiation-sector development \cite{Wells2026Radiation};
the later Planck/gravity analysis treats that bridge as a separate
theorem path rather than as part of the alpha fixed-point theorem.

This Planck/action result is relevant here only to empirical route
correspondence, because several operational extractions of \(\alpha\)
combine electromagnetic coupling, action, mass, and frequency/length
readouts. The separation of layers is simple: the alpha theorem derives
\(\alpha_F\); the Planck/action theorem supplies separate structural
action content; and empirical route analysis describes how operational
determinations combine electromagnetic coupling, action, mass, charge,
and readout conventions.

This paper uses the Planck/action bridge only as context for empirical
route correspondence. It does not merge the Planck derivation into the
alpha fixed-point theorem.

\subsection{Numerical agreement as consistency check}

The compact-fiber structural prediction is
\[
\alpha_F^{-1}=x_{\rm fiber}.
\]

The numerical agreement with the empirical low-energy fine-structure constant supports the role identification
\[
\alpha_F=\alpha_{\rm low-energy}.
\]

However, the numerical agreement is not the sole basis for the identification. The primary basis is structural: \(\alpha_F\) is the material/radiative bridge coupling, and leading atomic binding reads
\[
\alpha_F^2.
\]

The empirical comparison therefore functions as a consistency check on the role identification, not as the premise from which the identification is made.
\section{Correspondence with Empirical Alpha Routes}

The compact-fiber derivation gives a structural electromagnetic bridge
coupling
\[
\alpha_F=\frac1{x_{\rm fiber}},
\]
where \(x_{\rm fiber}\) is the positive root of the compact-fiber bridge
polynomial. The physical identification
\[
\alpha_F=\alpha_{\rm low-energy}
\]
rests on the role of this coupling in atomic binding: leading
material-radiative binding contains two bridge vertices and therefore
carries the scale
\[
\alpha_F^2.
\]

Empirical determinations of the fine-structure constant do not all
measure a single primitive object directly. They infer the low-energy
electromagnetic coupling through different physical systems: atomic
binding, spectroscopy, atom recoil, magnetic moments, and electrical
metrology
\cite{NIST_CODATA_2022_alpha_inv,Parker2018_alpha_Cs,
Morel2020_alpha_Rb,Aoyama2019_electron_gminus2_review,
BIPM_SI_Brochure_9th}. The compact-fiber result should therefore be
compared to empirical alpha science at the level of route structure. The
fixed-point theorem supplies the structural coupling \(\alpha_F\); the
empirical routes read that coupling through route-specific dimensional
anchors, correction schemes, and operational conventions. The general
compact-fiber readout and matching-correction layer is treated separately
in \cite{Wells2026FiniteReadout}.

\subsection{Atomic binding and spectroscopy}

The most direct correspondence is the atomic binding route. In
conventional atomic physics, the Hartree energy is
\[
E_h=\alpha^2m_ec^2,
\]
and the Rydberg constant is
\[
R_\infty=\frac{\alpha^2m_ec}{2h}.
\]
Both formulas place the squared electromagnetic coupling at the leading
Coulomb scale:
\[
\text{atomic binding scale}\propto \alpha^2.
\]

The compact-fiber source-response structure gives the corresponding
relation
\[
E_{\rm atomic,F}
=
C_{\rm mass/readout}\alpha_F^2,
\]
where \(C_{\rm mass/readout}\) contains dimensional mass, action, and
finite-readout matching factors \cite{Wells2026FiniteReadout}. Thus the
leading-order correspondence is
\[
\alpha^2
\quad
\leftrightarrow
\quad
\alpha_F^2.
\]

This makes the Hartree--Rydberg route the strongest empirical bridge for
the present result. The same squared-coupling role appears on both sides.
What remains external to the alpha fixed-point theorem is not the
electromagnetic coupling itself, but the dimensional anchoring and
correction stack: mass readout, action readout, reduced-mass corrections,
nuclear-size corrections, and the conversion from structural energy scale
to conventional frequency or length units.

Spectroscopic routes have the same leading structure. For hydrogenic
systems, the gross scale is controlled by expressions equivalent to
\[
E_n\sim -\frac{\mu c^2\alpha^2}{2n^2},
\]
with additional contributions from fine structure, Lamb shifts, hyperfine
structure, finite nuclear size, recoil, and system-specific corrections.
The compact-fiber correspondence is therefore strongest at the leading
Coulomb scale, where spectroscopy reads \(\alpha_F^2\) through the same
two-vertex material-radiative source-response structure. Higher-precision
spectroscopic extraction is a correction-route problem, not part of the
alpha fixed-point theorem itself.

\subsection{\texorpdfstring{Recoil and \(h/m\)-based determinations}{Recoil and h/m-based determinations}}

Atom recoil determinations infer \(\alpha\) through relations involving
recoil measurements, mass ratios, the Rydberg constant, and the action
anchor
\[
h.
\]
These routes include high-precision cesium and rubidium recoil
determinations and reviews of the \(h/m\) route to \(\alpha\)
\cite{Parker2018_alpha_Cs,Morel2020_alpha_Rb,
Clade2019_alpha_h_over_m_review}. Schematically, such routes combine
\[
R_\infty,
\qquad
h/m,
\qquad
m_e,
\qquad
c,
\]
together with measured mass ratios and recoil observables. They are
therefore not measurements of a bare coupling alone; they infer the
coupling through an action/mass/readout chain.

In compact-fiber terms, the route decomposes into two structural layers.
The alpha fixed-point theorem supplies
\[
\alpha_F,
\]
while the compact-fiber Planck/action and gravitational-coupling analysis
supplies separate structural content for the action anchor
\[
h
\]
\cite{Wells2026PlanckGravity}. In the notation of that work, the
Planck/action calculation uses the effective structural factor
\[
R_{\rm eff}
=
\frac{R_1}{(1+m+e)^{1/3}},
\]
together with the Larson/Nehru natural-unit anchor system.

Thus recoil and \(h/m\)-based routes now involve two distinct
compact-fiber theorem chains: the electromagnetic bridge coupling
\(\alpha_F\) from the present paper, and the Planck/action bridge from
the separate Planck/gravity analysis. This does not close the recoil
route. A complete compact-fiber treatment would still require structural
mass ratios, recoil kinematics, reduced-mass corrections, and
route-specific metrological conventions. The significance is narrower:
these routes are no longer merely single-input correspondences to
\(\alpha_F\); they provide a setting in which separate compact-fiber
structures for electromagnetic coupling and action enter the same
empirical extraction pipeline.

\subsection{Magnetic-moment and electrical-metrology routes}

The electron anomalous magnetic moment route extracts \(\alpha\) through
the QED radiative expansion for
\[
a_e=\frac{g-2}{2}.
\]
In standard form, this expansion is expressed in powers of
\[
\frac{\alpha}{\pi},
\]
with higher-order hadronic, electroweak, and mass-dependent
contributions entering at the precision frontier
\cite{Aoyama2019_electron_gminus2_review,Aoyama2012_tenth_order_QED}.

The compact-fiber correspondence is that \(g-2\) is an amplitude-level
radiative readout of the same low-energy electromagnetic coupling. The
present paper does not derive the QED vertex expansion, the coefficient
sequence, or the amplitude-level radiative normalization. A compact-fiber
derivation of this route would require an amplitude-level radiative
bridge formalism, a compact-fiber interpretation of vertex corrections,
and a derivation or reconstruction of the \(\alpha/\pi\) expansion
structure. In the present paper, the \(g-2\) route is treated as a
correspondence to the same coupling rather than as an independently
derived measurement route.

Electrical-standard routes involve relations such as
\[
K_J=\frac{2e}{h}
\]
for the Josephson constant and
\[
R_K=\frac{h}{e^2}
\]
for the von Klitzing constant. These routes connect charge, action,
impedance, and electrical metrology to the fine-structure constant
\cite{BIPM_SI_Brochure_9th}.

The alpha theorem supplies the electromagnetic bridge coupling
\(\alpha_F\), and the Planck/action work supplies structural action
content for \(h\) \cite{Wells2026PlanckGravity}. Electrical-standard
determinations also require a compact-fiber account of charge readout,
impedance, and metrological convention. They therefore provide a partial
correspondence at present: they involve the same electromagnetic
coupling, but a complete compact-fiber derivation of the route belongs to
a broader charge/action/electrical-metrology program.

\subsection{Structural coverage by empirical route}

The empirical-route correspondence can be summarized by structural
coverage rather than by route names alone. The table records where
compact-fiber structural content is presently available and which route
layers remain imported or open.

\begin{center}
{\footnotesize
\setlength{\tabcolsep}{2pt}
\begin{tabularx}{\textwidth}{
  >{\hsize=0.65\hsize\raggedright\arraybackslash\hspace{0pt}}X
  >{\hsize=0.95\hsize\raggedright\arraybackslash\hspace{0pt}}X
  >{\hsize=1.25\hsize\raggedright\arraybackslash\hspace{0pt}}X
  >{\hsize=1.55\hsize\raggedright\arraybackslash\hspace{0pt}}X
  >{\hsize=0.60\hsize\raggedright\arraybackslash\hspace{0pt}}X}
\toprule
Route & Quantities used & Compact-fiber structural content & Remaining imports or open layers & Coverage \\
\midrule
Hartree--Rydberg & \(\alpha^2\), \(m_e\), \(c\) & \(\alpha_F\) from the present bridge equation & electron-mass readout, dimensional/unit conventions & coupling \\
Spectroscopy & \(\alpha^2\) plus transition corrections & \(\alpha_F\) at the leading Coulomb scale & fine structure, Lamb shift, hyperfine, recoil, finite-size, and route-specific corrections & coupling \\
Recoil / \(h/m\) & \(\alpha\), \(R_\infty\), \(h/m_X\), mass ratios & \(\alpha_F\) from this paper; Planck/action content & mass ratios, recoil kinematics, reduced-mass and route-specific correction pipeline & coupling + action \\
Electron \(g{-}2\) & QED expansion in \(\alpha/\pi\) & \(\alpha_F\) as the low-energy coupling entering the amplitude expansion & QED vertex expansion, coefficient sequence, hadronic/electroweak terms, amplitude-level compact-fiber readout & coupling \\
Electrical standards & \(e\), \(h\), \(R_K\), \(K_J\), impedance relations & \(\alpha_F\); Planck/action content & charge readout, impedance normalization, electrical-metrology conventions & coupling + action \\
\bottomrule
\end{tabularx}
}
\end{center}

This table is not a list of completed route derivations. It records where
the compact-fiber structural coupling and related compact-fiber theorem
paths enter existing empirical alpha science.

\subsection{Multi-route consistency}

The compact-fiber bridge coupling \(\alpha_F\) is derived without using
any route-specific empirical extraction of \(\alpha\) as an algebraic
input. Hartree/Rydberg spectroscopy, atom-recoil routes, electron
\(g{-}2\), and electrical-standard routes access the low-energy
electromagnetic coupling through different auxiliary structures:
electron mass, action, recoil kinematics, radiative vertex coefficients,
charge, and impedance conventions.

This route diversity strengthens the role identification because the
routes place the electromagnetic coupling in different auxiliary
structures: atomic binding, recoil/action chains, radiative vertex
expansions, and electrical metrology. It should not, however, be read as
a completed metrological closure. The present paper does not derive every
route-specific correction pipeline, nor does it assign a combined
statistical likelihood over the empirical route network. The point is
more limited: the same compact-fiber bridge coupling occupies the
coupling slot that these different operational routes infer. Agreement
with the empirical low-energy value is therefore a consistency check of
the role identification
\[
\alpha_F=\alpha_{\rm low-energy},
\]
not merely a comparison to one favored extraction route.

The important point is that the compact-fiber value is not identified
with the fine-structure constant merely because \(x_{\rm fiber}\) is
numerically close to the measured inverse value. The identification rests
on the structural role of \(\alpha_F\) as the material/radiative bridge
coupling and on the fact that the leading atomic scale reads that bridge
through \(\alpha_F^2\). The major empirical routes then appear as
different operational readouts of the same low-energy coupling, each with
its own dimensional anchors and correction structure.
\section{Scope, Limitations, and Falsifiability}

The compact-fiber derivation gives a structural electromagnetic bridge
coupling
\[
\alpha_F=\frac{1}{x_{\rm fiber}},
\]
where \(x_{\rm fiber}\) is the positive root of
\[
x^4
-
137x^3
-
\frac{74}{15}x^2
+
\frac{296}{105}
=
0.
\]
The claim is conditional: given the compact-fiber carrier, the
radiative/vibrational bridge-placement readout, the carrier-access
composition rule, the closure-budget accounting, the self-consistent
bridge-sharing rule, and the closed-return correction stated above, the
fixed-point equation and its positive root follow algebraically.

The result should therefore be read as a conditional structural
derivation of a low-energy electromagnetic bridge coupling, not as an
unconditional derivation from bare postulates alone and not as a completed
derivation of every empirical route by which \(\alpha\) is extracted.
Empirical determinations of the fine-structure constant involve
additional anchors and correction structures. For example,
\[
R_\infty=\frac{\alpha^2m_ec}{2h}
\]
contains the electromagnetic coupling, but also mass, action, and
unit-readout structure. Recoil, \(g-2\), spectroscopy, Josephson, and
quantum Hall determinations likewise infer \(\alpha\) through
route-specific theoretical and metrological pipelines.

The present paper separates these layers. The compact-fiber fixed-point
equation gives the structural bridge coupling \(\alpha_F\). The physical
role identification
\[
\alpha_F=\alpha_{\rm low-energy}
\]
rests on the \(\alpha_F^2\) role in leading atomic binding.
Route-specific metrological closure remains downstream; the general
compact-fiber readout and matching-correction layer is treated in
\cite{Wells2026FiniteReadout}.

\subsection{Coefficient status}

The fixed-point equation contains no continuously adjusted coefficient
and uses no measured value of \(\alpha\) as an algebraic input. Its
coefficients have named compact-fiber origins:
\[
137=128+9,
\qquad
\frac{74}{15}=5-\frac1{15},
\qquad
\frac{296}{105}=\left(\frac{74}{15}\right)\left(\frac47\right).
\]
Here
\[
128=|Z_4\times Z_4\times Z_8|,
\qquad
9=|V_{\rm rad}^{\rm bridge}|=3^2,
\]
\[
5=N_m+1,
\qquad
\frac1{15}=\frac1{N_m^2-1},
\qquad
\frac47=\frac{N_m}{N_e-1}.
\]
Thus the polynomial has a fully specified compact-fiber coefficient
structure under the stated bridge-readout and closure rules. The
empirical value of \(\alpha\) is not used to set the base count, the
feedback burden, or the closed-return coefficient.

This distinction is important. The paper does not claim that every
bridge-readout and closure rule used here is derived from bare postulates
within this manuscript. It claims that, once those rules are stated, the
alpha bridge equation follows without continuously adjusted coefficients
and without using the measured value of \(\alpha\) to set the algebraic
coefficients.

\subsection{Claim boundary}

The derivation depends on the following load-bearing ingredients:
\begin{enumerate}[label=(\arabic*)]
\item the compact-fiber carrier
\[
F=Z_4\times Z_4\times Z_8;
\]

\item the sign-status and placement-readout construction
\[
|V_{\rm rad}^{\rm bridge}|=3^2=9;
\]

\item carrier-level disjoint-union access counting for the
material/radiative bridge;

\item closure-budget accounting for the feedback burden
\[
K=\frac{74}{15};
\]

\item self-consistent transport sharing,
\[
x=B+\frac{K}{x};
\]

\item the source/complement return coefficient
\[
\frac{N_m}{N_e-1};
\]

\item the derivative self-response sign of the closed return.
\end{enumerate}

The result would require revision if any of these ingredients failed. In
particular, the derivation would require revision if \(F\) were not the
relevant material closure carrier, if
\(|V_{\rm rad}^{\rm bridge}|=3^2=9\) were not the correct
bridge-placement count, if the electromagnetic bridge required
product-state counting \(F\times V_{\rm rad}^{\rm bridge}\) rather than
disjoint-union access, or if the feedback burden \(K=74/15\) were not a
valid closure-budget term. It would also require revision if the
second-order return term
\[
-\frac{K N_m}{(N_e-1)x^3}
\]
were not the correct closed-return correction, or if another same-order
scalar bridge correction were admitted.

Finally, the low-energy role identification would fail if \(\alpha_F\)
did not play the \(\alpha^2\) role in leading atomic binding, or if
future measurement-chain analysis showed that operational extractions of
\(\alpha\) read a systematically different quantity than the
compact-fiber bridge coupling.

\subsection{Same-order and higher-order boundaries}

The second-order correction is conditionally unique within the current
scalar two-node bridge graph. A same-order competitor would require at
least one of the following:
\begin{enumerate}[label=(\arabic*)]
\item a second independent first-order feedback source;
\item a new carrier pool;
\item a preimage-resolved observable rather than scalar bridge readout;
\item a radiation-side feedback mechanism contributing to \(K\);
\item an independent third transport node;
\item a direct cubic mechanism not generated by the first-order feedback.
\end{enumerate}

No such mechanism is admitted in the present derivation. The uniqueness
claim is therefore relative to the stated scalar bridge graph and the
framework rules used in this paper. It is not a claim that all future
radiation-side feedback mechanisms are impossible.

The quartic equation includes the first closed return of the first-order
feedback burden. If repeated closed returns are admitted, each additional
round trip is suppressed by at least
\[
\frac{C}{x^2},
\qquad
C=\frac{N_m}{N_e-1}=\frac{4}{7}.
\]
For
\[
x\approx137.036,
\]
this suppression factor is approximately
\[
3.0\times10^{-5}.
\]
Thus higher-order repeated returns, if present, are operationally small
at the current comparison scale. A complete all-orders return theorem
remains open.

\subsection{Numerical comparison and route scope}

The compact-fiber point value lies within the CODATA 2022 experimental
uncertainty band for the inverse fine-structure constant. This comparison
is not a statistical validation of the framework, since the present paper
does not assign a theoretical uncertainty to the imported bridge-readout
rules, the closure-budget rules, the closed-return construction, or
possible omitted same-order mechanisms. It should therefore be read as a
consistency check, not as an error-budget closure.

Some empirical routes to \(\alpha\), especially Rydberg/recoil and
\(h/m\)-based routes, involve the action anchor
\[
h.
\]
The compact-fiber program treats the structural content of \(h\) as a
separate theorem path \cite{Wells2026PlanckGravity}. That path is
relevant to empirical route correspondence, but it is not part of the
alpha fixed-point theorem. The
alpha theorem supplies the compact-fiber bridge coupling \(\alpha_F\);
the Planck/action theorem supplies separate structural action content;
route analysis then explains how empirical determinations combine
electromagnetic coupling, action, mass, charge, and readout conventions.

The Planck/action result is used here only as context for legacy-route
correspondence.

\subsection{Falsifiability and summary of scope}

The present paper supplies an internal structural prediction and a
comparison with the empirical low-energy coupling. It does not claim a
complete independent falsification program. The internal derivation is
falsifiable at the framework level by failure of any stated load-bearing
rule: carrier choice, bridge-placement count, disjoint-union access,
closure-budget accounting, self-consistent sharing, return coefficient,
or derivative self-response sign.

A stronger external test would require applying the same bridge-readout
and closure rules to a second observable without introducing new
structural choices. Such a test would distinguish a reusable compact-fiber
rule set from an isolated alpha-specific construction. The present paper
identifies the alpha bridge prediction and its failure modes, while
leaving broader route-by-route and second-observable closure to
subsequent work.

Within the stated framework boundaries, the compact-fiber derivation gives
\[
\boxed{
\alpha_F^{-1}=x_{\rm fiber}.
}
\]
The result is not merely numerical: each coefficient in the fixed-point
equation has a named structural origin. It is also not a completed
metrological derivation: route-specific empirical extraction pipelines
remain downstream theorem programs. The strongest claim of this paper is
therefore a conditional structural prediction for the low-energy
electromagnetic bridge coupling, with explicit route-readout,
theoretical-uncertainty, and independent-test boundaries.
\section{Conclusion}

The compact-fiber carrier $F = \mathbb{Z}_4 \times \mathbb{Z}_4
\times \mathbb{Z}_8$ and the radiative/vibrational
bridge-placement carrier $V_{\rm rad}^{\rm bridge} \cong
\{-,0,+\}^2$ supply the base electromagnetic bridge count
$B = |F| + |V_{\rm rad}^{\rm bridge}| = 128 + 9 = 137$ through
mechanism-distinct disjoint-union access.  Cover-mediated
feedback with closure-budget accounting gives the net burden
$K = (N_m{+}1) - 1/(N_m^2{-}1) = 74/15$, and self-consistent
transport sharing produces the first-order bridge equation
$x = B + K/x$.  The first closed return on the scalar two-node
bridge graph contributes
$-K N_m/\bigl((N_e{-}1)\,x^3\bigr)$, yielding the quartic
\[
  x^4 - 137\,x^3 - \tfrac{74}{15}\,x^2
  + \tfrac{296}{105} = 0,
\]
whose positive root is
$x_{\rm fiber} = 137.0359991771859\ldots$, so
$\alpha_F = 1/x_{\rm fiber}$.  No measured value of~$\alpha$
is used as an algebraic input, and no continuously adjusted
coefficient appears in the equation.

The identification $\alpha_F = \alpha_{\rm low\text{-}energy}$
is made by role: the compact-fiber electromagnetic bridge is the
material/radiative interface, and leading atomic binding reads
that bridge through two vertices, giving
$E_{\rm atomic} \propto \alpha_F^2$.  This matches the role
of~$\alpha^2$ in the Hartree/Rydberg scale and provides the
structural basis for identifying~$\alpha_F$ with the low-energy
fine-structure constant.

The numerical comparison with the CODATA~2022 value
$\alpha^{-1} = 137.035999177(21)$ is a consistency check: the
compact-fiber point value lies within the experimental
uncertainty band.  This comparison is not an error-budget
closure, since the paper assigns no theoretical uncertainty
to the imported bridge-readout rules or to possible omitted
same-order mechanisms.

The result is conditional on the named compact-fiber framework
lemmas: carrier choice, bridge-placement count, disjoint-union
access, closure-budget accounting, self-consistent sharing,
return coefficient, and derivative self-response sign.  It is
not an unconditional derivation from bare postulates alone, and
it is not a completed route-by-route derivation of CODATA or
metrological extraction pipelines.  Route-specific closure and
independent second-observable testing remain downstream work.

\bibliographystyle{plain}
\bibliography{sections/references}

\end{document}
